% GENERATED FILE. Edit canonical lessons, then run npm run generate:books.
% canonical-source-hash: fnv1a64:d78ba0a7abe11d56
% canonical-lessons: ML-C397-rajavu, ML-C397-rajni, ML-C397-rajakumaran, ML-C397-rajakumari, ML-C397-kottaram

\chapter{Kings and Queens}
\label{ch:ml-kings-and-queens}
\hlchaptermodality{397}

\begin{chapteropening}
\textbf{By the end of this chapter:} \emph{I can say a king, a queen, a prince, a princess and a palace.}

Complete the last lesson of chapter 397: I can say a king, a queen, a prince, a princess and a palace.
\end{chapteropening}

\section[\texorpdfstring{\ml{രാജാവ്} (rājāvŭ)}{രാജാവ് (rājāvŭ)}]{\ml{രാജാവ്} (rājāvŭ) — a king}

\label{lesson:ML-C397-rajavu}



\begin{quote}
Before the new one: say the Malayalam for to honour, then the Malayalam for to make one's own.
\end{quote}

\subsection*{You'll want to know: \ml{രാജാവ്}}
\textbf{\ml{രാജാവ്}} — \emph{rājāvŭ} — \textquotedblleft{}a king\textquotedblright{}.

\begin{quote}
\ml{പണ്ടുപണ്ട്} \ml{ഒരു} \ml{രാജ്യത്ത്} \ml{ഒരു} \ml{നല്ല} \ml{രാജാവ്} \ml{ഉണ്ടായിരുന്നു}.
\end{quote}

\begin{quote}
\emph{paṇṭupaṇṭŭ oru rājyattŭ oru nalla rājāvŭ uṇṭāyirunnu.}
\end{quote}

\begin{quote}
Once upon a time there was a good king in a country.
\end{quote}

Plural: \ml{രാജാക്കന്മാർ}, \emph{rājākkanmār}.

Plural \ml{രാജാക്കന്മാർ}.

\begin{cousinweb}
\textbf{\ml{രാജാവ്}} (\emph{rājāvŭ}) is borrowed from Sanskrit \emph{rājan}, \textquotedblleft{}a chieftain, a king\textquotedblright{} (CDIAL 10679), the root of \ml{രാജ്യം} (\emph{rājyaṁ}), \textquotedblleft{}a country, a kingdom\textquotedblright{}.
\end{cousinweb}

\begin{grammarlens}[title={What you've built}]
One more word for telling what happened.
\end{grammarlens}

\subsection*{Your turn}
\begin{itemize}
\raggedright
  \item \emph{Say it:} \emph{rājāvŭ}
  \item \emph{Say it:} \emph{rājāvŭ}, once more
  \item \emph{Say it:} \emph{paṇṭupaṇṭŭ oru rājyattŭ oru nalla rājāvŭ uṇṭāyirunnu.} — Once upon a time there was a good king in a country.
  \item \emph{Recall:} say the Malayalam for to honour, then the Malayalam for to make one's own, then say \emph{rājāvŭ} again
\end{itemize}

\begin{tcolorbox}[breakable,colback=teal!4,colframe=teal!35!black,title={Before you move on}]
What does \ml{രാജാവ്} mean? (\textquotedblleft{}A king\textquotedblright{}.) Say it once more.
\end{tcolorbox}
\section[\texorpdfstring{\ml{രാജ്ഞി} (rājñi)}{രാജ്ഞി (rājñi)}]{\ml{രാജ്ഞി} (rājñi) — a queen}

\label{lesson:ML-C397-rajni}



\begin{quote}
Before the new one: say the Malayalam for to make one's own, then the Malayalam for a king.
\end{quote}

\subsection*{You'll want to know: \ml{രാജ്ഞി}}
\textbf{\ml{രാജ്ഞി}} — \emph{rājñi} — \textquotedblleft{}a queen\textquotedblright{}.

\begin{quote}
\ml{രാജാവും} \ml{രാജ്ഞിയും} \ml{ഗ്രാമം} \ml{സന്ദർശിച്ചു}.
\end{quote}

\begin{quote}
\emph{rājāvuṁ rājñiyuṁ grāmaṁ sandarśiccu.}
\end{quote}

\begin{quote}
The king and the queen visited the village.
\end{quote}

The queen who rules in her own right, or the king's wife.

\begin{cousinweb}
\textbf{\ml{രാജ്ഞി}} (\emph{rājñi}) is borrowed from Sanskrit \emph{rājñī}, \textquotedblleft{}a queen\textquotedblright{} (CDIAL 10692), the feminine of \emph{rājan}.
\end{cousinweb}

\begin{grammarlens}[title={What you've built}]
One more word for telling what happened.
\end{grammarlens}

\subsection*{Your turn}
\begin{itemize}
\raggedright
  \item \emph{Say it:} \emph{rājñi}
  \item \emph{Say it:} \emph{rājñi}, once more
  \item \emph{Say it:} \emph{rājāvuṁ rājñiyuṁ grāmaṁ sandarśiccu.} — The king and the queen visited the village.
  \item \emph{Recall:} say the Malayalam for to make one's own, then the Malayalam for a king, then say \emph{rājñi} again
\end{itemize}

\begin{tcolorbox}[breakable,colback=teal!4,colframe=teal!35!black,title={Before you move on}]
What does \ml{രാജ്ഞി} mean? (\textquotedblleft{}A queen\textquotedblright{}.) Say it once more.
\end{tcolorbox}
\section[\texorpdfstring{\ml{രാജകുമാരൻ} (rājakumāran)}{രാജകുമാരൻ (rājakumāran)}]{\ml{രാജകുമാരൻ} (rājakumāran) — a prince}

\label{lesson:ML-C397-rajakumaran}



\begin{quote}
Before the new one: say the Malayalam for a king, then the Malayalam for a queen.
\end{quote}

\subsection*{You'll want to know: \ml{രാജകുമാരൻ}}
\textbf{\ml{രാജകുമാരൻ}} — \emph{rājakumāran} — \textquotedblleft{}a prince\textquotedblright{}.

\begin{quote}
\ml{രാജകുമാരൻ} \ml{കുതിരപ്പുറത്ത്} \ml{കാട്ടിലേക്ക്} \ml{പോയി}.
\end{quote}

\begin{quote}
\emph{rājakumāran kutirappuṟattŭ kāṭṭilēkkŭ pōyi.}
\end{quote}

\begin{quote}
The prince rode into the forest on horseback.
\end{quote}

\ml{രാജ}-, king, + \ml{കുമാരൻ}, a boy: the king's son.

\begin{cousinweb}
\textbf{\ml{രാജകുമാരൻ}} (\emph{rājakumāran}) is built from Sanskrit \emph{rāja-}, \textquotedblleft{}king\textquotedblright{} (CDIAL 10679), and \emph{kumāra}, \textquotedblleft{}a boy; a prince\textquotedblright{} (CDIAL 3303).
\end{cousinweb}

\begin{grammarlens}[title={What you've built}]
One more word for telling what happened.
\end{grammarlens}

\subsection*{Your turn}
\begin{itemize}
\raggedright
  \item \emph{Say it:} \emph{rājakumāran}
  \item \emph{Say it:} \emph{rājakumāran}, once more
  \item \emph{Say it:} \emph{rājakumāran kutirappuṟattŭ kāṭṭilēkkŭ pōyi.} — The prince rode into the forest on horseback.
  \item \emph{Recall:} say the Malayalam for a king, then the Malayalam for a queen, then say \emph{rājakumāran} again
\end{itemize}

\begin{tcolorbox}[breakable,colback=teal!4,colframe=teal!35!black,title={Before you move on}]
What does \ml{രാജകുമാരൻ} mean? (\textquotedblleft{}A prince\textquotedblright{}.) Say it once more.
\end{tcolorbox}
\section[\texorpdfstring{\ml{രാജകുമാരി} (rājakumāri)}{രാജകുമാരി (rājakumāri)}]{\ml{രാജകുമാരി} (rājakumāri) — a princess}

\label{lesson:ML-C397-rajakumari}



\begin{quote}
Before the new one: say the Malayalam for a queen, then the Malayalam for a prince.
\end{quote}

\subsection*{You'll want to know: \ml{രാജകുമാരി}}
\textbf{\ml{രാജകുമാരി}} — \emph{rājakumāri} — \textquotedblleft{}a princess\textquotedblright{}.

\begin{quote}
\ml{രാജകുമാരി} \ml{പൂന്തോട്ടത്തിൽ} \ml{പാട്ട്} \ml{പാടി}.
\end{quote}

\begin{quote}
\emph{rājakumāri pūntōṭṭattil pāṭṭŭ pāṭi.}
\end{quote}

\begin{quote}
The princess sang in the garden.
\end{quote}

The feminine of \ml{രാജകുമാരൻ}.

\begin{cousinweb}
\textbf{\ml{രാജകുമാരി}} (\emph{rājakumāri}) is built from Sanskrit \emph{rāja-}, \textquotedblleft{}king\textquotedblright{} (CDIAL 10679), and \emph{kumārī}, \textquotedblleft{}a girl, a princess\textquotedblright{} (CDIAL 3303).
\end{cousinweb}

\begin{grammarlens}[title={What you've built}]
One more word for telling what happened.
\end{grammarlens}

\subsection*{Your turn}
\begin{itemize}
\raggedright
  \item \emph{Say it:} \emph{rājakumāri}
  \item \emph{Say it:} \emph{rājakumāri}, once more
  \item \emph{Say it:} \emph{rājakumāri pūntōṭṭattil pāṭṭŭ pāṭi.} — The princess sang in the garden.
  \item \emph{Recall:} say the Malayalam for a queen, then the Malayalam for a prince, then say \emph{rājakumāri} again
\end{itemize}

\begin{tcolorbox}[breakable,colback=teal!4,colframe=teal!35!black,title={Before you move on}]
What does \ml{രാജകുമാരി} mean? (\textquotedblleft{}A princess\textquotedblright{}.) Say it once more.
\end{tcolorbox}
\section[\texorpdfstring{\ml{കൊട്ടാരം} (koṭṭāraṁ)}{കൊട്ടാരം (koṭṭāraṁ)}]{\ml{കൊട്ടാരം} (koṭṭāraṁ) — a palace}

\label{lesson:ML-C397-kottaram}



\begin{quote}
Before the new one: say the Malayalam for a prince, then the Malayalam for a princess.
\end{quote}

\subsection*{You'll want to know: \ml{കൊട്ടാരം}}
\textbf{\ml{കൊട്ടാരം}} — \emph{koṭṭāraṁ} — \textquotedblleft{}a palace\textquotedblright{}.

\begin{quote}
\ml{രാജാവിന്റെ} \ml{കൊട്ടാരം} \ml{കുന്നിന്റെ} \ml{മുകളിൽ} \ml{ആയിരുന്നു}.
\end{quote}

\begin{quote}
\emph{rājāvinṟe koṭṭāraṁ kunninṟe mukaḷil āyirunnu.}
\end{quote}

\begin{quote}
The king's palace was on top of the hill.
\end{quote}

Also any grand old mansion.

\begin{cousinweb}
\textbf{\ml{കൊട്ടാരം}} (\emph{koṭṭāraṁ}) is usually traced to Sanskrit \emph{koṣṭhāgāra}, \textquotedblleft{}a store-room, a treasury\textquotedblright{} (Monier-Williams): the king's house was where his stores were kept.
\end{cousinweb}

\begin{grammarlens}[title={What you've built}]
One more word for telling what happened.
\end{grammarlens}

\subsection*{Your turn}
\begin{itemize}
\raggedright
  \item \emph{Say it:} \emph{koṭṭāraṁ}
  \item \emph{Say it:} \emph{koṭṭāraṁ}, once more
  \item \emph{Say it:} \emph{rājāvinṟe koṭṭāraṁ kunninṟe mukaḷil āyirunnu.} — The king's palace was on top of the hill.
  \item \emph{Recall:} say the Malayalam for a prince, then the Malayalam for a princess, then say \emph{koṭṭāraṁ} again
\end{itemize}

\begin{tcolorbox}[breakable,colback=teal!4,colframe=teal!35!black,title={Before you move on}]
What does \ml{കൊട്ടാരം} mean? (\textquotedblleft{}A palace\textquotedblright{}.) Say it once more.
\end{tcolorbox}
